% ============================================================================
%  Laboratorio di Programmazione -- Lezione 9 (3 novembre 2026)
%  Compilare con:  pdflatex -shell-escape Lezione09_Interfacce_Package.tex
%  Richiede: mermaid.sty (nel progetto) + Mermaid CLI `mmdc` nel PATH
%  NB: ogni frame che contiene codice (lstlisting) o un diagramma (mermaid)
%      deve essere dichiarato [fragile].
% ============================================================================
\documentclass[10pt,aspectratio=169]{beamer}
\usetheme{Madrid}
\usecolortheme{default}

% `provide=*` carica la localizzazione italiana dai file .ini di babel: funziona anche
% su installazioni TeX prive di texlive-lang-italian; con una TeX Live completa e`
% equivalente al classico \usepackage[italian]{babel}.
\usepackage[italian,provide=*]{babel}
\usepackage[utf8]{inputenc}

% Localizzazione dei nomi dei blocchi Beamer
\uselanguage{italian}
\languagepath{italian}
\deftranslation[to=italian]{Definition}{Definizione}
\deftranslation[to=italian]{Theorem}{Teorema}
\deftranslation[to=italian]{Example}{Esempio}
\deftranslation[to=italian]{Proof}{Dimostrazione}
\deftranslation[to=italian]{Corollary}{Corollario}
\deftranslation[to=italian]{Lemma}{Lemma}

\usepackage{listings}
\usepackage{xcolor}
\usepackage{hyperref}
\usepackage{array}
\usepackage{booktabs}
\usepackage{mermaid}

% Mermaid: sfondo trasparente e dimensione massima adatta a una slide 16:9
\MermaidRendererOptions{-b transparent}
\MermaidAdjustBoxOpts{max width=0.95\linewidth,max height=0.62\textheight,center}

% Configurazione colori per il codice
\definecolor{codegreen}{rgb}{0,0.6,0}
\definecolor{codegray}{rgb}{0.5,0.5,0.5}
\definecolor{codepurple}{rgb}{0.58,0,0.82}
\definecolor{backcolour}{rgb}{0.95,0.95,0.92}
\definecolor{keywordblue}{rgb}{0.0,0.0,0.7}

% Stile per snippet Java
\lstdefinestyle{javastyle}{
    backgroundcolor=\color{backcolour},
    commentstyle=\color{codegreen}\itshape,
    keywordstyle=\color{keywordblue}\bfseries,
    numberstyle=\scriptsize\color{codegray},
    stringstyle=\color{codepurple},
    basicstyle=\ttfamily\footnotesize,
    breakatwhitespace=false,
    breaklines=true,
    captionpos=b,
    keepspaces=true,
    numbers=none,
    numbersep=5pt,
    showspaces=false,
    showstringspaces=false,
    showtabs=false,
    tabsize=2,
    aboveskip=0.6\medskipamount,
    belowskip=0.6\medskipamount,
    language=Java,
    morekeywords={default},
    literate={à}{{\`a}}1 {è}{{\`e}}1 {é}{{\'e}}1 {ì}{{\`i}}1
             {ò}{{\`o}}1 {ù}{{\`u}}1 {È}{{\`E}}1
}
\lstset{style=javastyle}

% Stile per righe di terminale
\lstdefinestyle{shell}{
    language=bash,
    backgroundcolor=\color{black!85},
    basicstyle=\ttfamily\footnotesize\color{white},
    keywordstyle=\color{white},
    stringstyle=\color{white},
    commentstyle=\color{green!70},
    numbers=none,
    frame=none,
    aboveskip=0.6\medskipamount,
    belowskip=0.6\medskipamount,
    showstringspaces=false
}

\title[Lezione 9 -- Interfacce e package]{Lezione 9 -- Interfacce e package}
\subtitle{\texttt{interface} e \texttt{implements}, più tipi per una classe, metodi \texttt{default} e \texttt{static}; \texttt{package}, \texttt{import} e visibilità}
\author[Mattia Fonisto]{Mattia Fonisto}
\institute[Lab.\ di Programmazione]{Laboratorio di Programmazione -- Ingegneria Informatica, II anno\\
           Università degli Studi di Napoli Federico II -- Sede di San Giovanni a Teduccio}
\date{3 novembre 2026}

\begin{document}

\frame{\titlepage}

% ========================================================================
\begin{frame}{Agenda della lezione (2h)}
\begin{enumerate}
    \item \textbf{Ripasso e Homework 8} \hfill (\emph{15 min})
    \begin{itemize}
        \item uno di voi alla lavagna con \texttt{Data}, \texttt{Appuntamento} o \texttt{Anagrafe}: lo commentiamo insieme
    \end{itemize}
    \item \textbf{Le interfacce} \hfill (\emph{25 min})
    \begin{itemize}
        \item \texttt{interface} e \texttt{implements}; l'interfaccia come tipo; più interfacce per una classe
        \item metodi \texttt{default} e \texttt{static}; interfacce o classi astratte; finalmente \texttt{Cloneable}
    \end{itemize}
    \item \textbf{I package} \hfill (\emph{20 min})
    \begin{itemize}
        \item \texttt{package}, nomi qualificati, cartelle; compilare ed eseguire; \texttt{import}
        \item visibilità di package e riepilogo dei modificatori
    \end{itemize}
    \item \textbf{Esercitazione} -- tutti insieme \hfill (\emph{50 min})
    \begin{itemize}
        \item \emph{Commentiamo insieme}: un simulatore di dispositivi con due interfacce
        \item \emph{Insieme, passo passo}: lo stesso simulatore riorganizzato in package
    \end{itemize}
    \item \textbf{Wrap-up e Homework 9} \hfill (\emph{10 min})
\end{enumerate}
\end{frame}

% ========================================================================
\section{Le interfacce}

\begin{frame}[fragile]{Interfacce: il problema}
\begin{columns}[T]
\column{0.58\textwidth}
\footnotesize
Una \texttt{Fattura} e un \texttt{Dipendente}: classi diverse, nessuna relazione \emph{è-un} fra loro. Eppure la contabilità li tratta allo stesso modo: ognuno \textbf{va pagato}, di ognuno serve l'\textbf{importo}. Vorremmo un array con entrambi e un metodo \texttt{totale} che sommi gli importi.

\vspace{0.1cm}
Con l'ereditarietà (lezione~7), una superclasse astratta \texttt{Pagabile} con \texttt{importo()} astratto? Due problemi:
\begin{itemize}\itemsep=2pt
    \item \texttt{Pagabile} non avrebbe \textbf{nulla} da dare in eredità: niente variabili di istanza, costruttori, metodi con il corpo; solo l'intestazione di \texttt{importo()}
    \item una classe ha \textbf{una sola} superclasse: se \texttt{Dipendente} estendesse già \texttt{Persona}, non potrebbe estendere \texttt{Pagabile}
\end{itemize}
\begin{block}{Interfaccia}
\small Un'\emph{interfaccia} (\texttt{interface}) è un tipo che elenca le intestazioni dei metodi che un oggetto offre, di norma senza corpo e sempre senza stato: dice \textbf{che cosa} si può chiedere a un oggetto, non \textbf{come} lo fa. Una classe che la \emph{implementa} si impegna a scrivere quei metodi. È l'interfaccia della lezione~4 come tipo di Java.
\end{block}

\column{0.42\textwidth}
\MermaidAdjustBoxOpts{max width=\linewidth,max height=0.52\textheight,center}
\begin{mermaid}
classDiagram
  direction TB
  class Pagabile {
    <<interface>>
    +importo() double
  }
  class Fattura {
    -numero : String
    -quantita : int
    -prezzo : double
    +importo() double
  }
  class Dipendente {
    -nome : String
    -stipendio : double
    +aumenta(double)
    +importo() double
  }
  Pagabile <|.. Fattura
  Pagabile <|.. Dipendente
\end{mermaid}
\MermaidAdjustBoxOpts{max width=0.95\linewidth,max height=0.62\textheight,center}
\footnotesize
In UML: \texttt{<<interface>>} sopra il nome; la freccia \textbf{tratteggiata} con il triangolo è la \emph{realizzazione}: \texttt{Fattura} implementa \texttt{Pagabile}.
\end{columns}
\end{frame}

% ========================================================================
\begin{frame}[fragile]{Dichiarare un'interfaccia -- \texttt{Pagabile.java}}
\begin{columns}[T]
\column{0.5\textwidth}
\begin{lstlisting}[basicstyle=\ttfamily\scriptsize]
/** Qualcosa che va pagato: ha un importo. */
public interface Pagabile {

    double importo();  // public abstract
}
\end{lstlisting}
\footnotesize
\begin{itemize}\itemsep=3pt
    \item \texttt{interface} al posto di \texttt{class}; i metodi hanno solo l'\textbf{intestazione} (la firma più il tipo di risultato), chiusa dal punto e virgola, come i metodi astratti della lezione~7
    \item sono \textbf{implicitamente} \texttt{public abstract}: scriverlo è lecito ma ridondante. Niente \texttt{protected}: ciò che un'interfaccia offre è tutto pubblico
    \item può dichiarare \textbf{costanti}: \texttt{double TARIFFA = 2.0;} è implicitamente \texttt{public static final}, si usa come \texttt{Pagabile.TARIFFA}
    \item \textbf{non} può avere variabili di istanza né costruttori: non descrive uno stato, solo un comportamento
    \item sta nel file \texttt{Pagabile.java}: ``una \texttt{public} per file, con il suo nome'' vale anche per le interfacce
\end{itemize}

\column{0.5\textwidth}
\begin{block}{Che cosa contiene un'interfaccia}
\small Intestazioni di metodi (astratti) e costanti. Da Java~8 anche metodi con il corpo, \texttt{default} e \texttt{static} (fra quattro slide). Mai variabili di istanza né costruttori.
\end{block}
\footnotesize
\begin{center}
\renewcommand{\arraystretch}{1.2}\setlength{\tabcolsep}{4pt}
\scriptsize
\begin{tabular}{@{}lll@{}}
\toprule
 & \textbf{classe} & \textbf{interfaccia} \\
\midrule
parola chiave & \texttt{class} & \texttt{interface} \\
variabili di istanza & sì & no \\
costanti \texttt{static final} & sì & sì (implicito) \\
costruttori & sì & no \\
metodi con il corpo & sì & \texttt{default} e \texttt{static} \\
metodi astratti & solo se \texttt{abstract} & sì, è la norma \\
\texttt{new} & sì (se concreta) & no \\
\bottomrule
\end{tabular}
\end{center}
\vspace{0.1cm}
Il nome è spesso un aggettivo: che cosa l'oggetto \emph{sa fare}: \texttt{Pagabile}, \texttt{Stampabile}, \texttt{Cloneable}, \texttt{Comparable}.
\end{columns}
\end{frame}

% ========================================================================
\begin{frame}[fragile]{Implementare un'interfaccia: \texttt{implements} -- \texttt{Fattura.java}}
\begin{columns}[T]
\column{0.51\textwidth}
\begin{lstlisting}[basicstyle=\ttfamily\scriptsize]
/** Una fattura: quantita` per prezzo. */
public class Fattura implements Pagabile {

    private final String numero;
    private final int quantita;
    private final double prezzo;

    public Fattura(String numero, int quantita,
                   double prezzo) {
        this.numero = numero;
        this.quantita = quantita;
        this.prezzo = prezzo;
    }

    @Override
    public double importo() {
        return quantita * prezzo;
    }
\end{lstlisting}
\footnotesize
Non mostrato il \texttt{toString} (\texttt{"fattura " + numero}). \texttt{Dipendente} è fatta allo stesso modo: \texttt{nome} e \texttt{stipendio}, costruttore, \texttt{importo()} che restituisce lo stipendio, \texttt{toString}, e in più \texttt{aumenta(double aumento)}.

\column{0.49\textwidth}
\begin{lstlisting}[basicstyle=\ttfamily\scriptsize]
/** Un dipendente: va pagato lo stipendio. */
public class Dipendente implements Pagabile {
\end{lstlisting}
\footnotesize
\begin{itemize}\itemsep=3pt
    \item \texttt{implements Pagabile} dopo il nome: la classe \textbf{promette} tutti i metodi dell'interfaccia. Se ne manca uno: \emph{``Fattura is not abstract and does not override abstract method importo() in Pagabile''}: o lo si scrive, o la classe diventa \texttt{abstract}
    \item il metodo va dichiarato \textbf{\texttt{public}}: nell'interfaccia lo è; con meno visibilità, \emph{``attempting to assign weaker access privileges''}. \texttt{@Override} vale anche qui
    \item il resto è libero: ogni classe ha le \emph{sue} variabili di istanza, i suoi costruttori e altri metodi, come \texttt{aumenta}, che \texttt{Pagabile} non conosce
\end{itemize}
\end{columns}
\end{frame}

% ========================================================================
\begin{frame}[fragile]{L'interfaccia come tipo: polimorfismo -- \texttt{ProvaPagabili.java}}
\begin{columns}[T]
\column{0.56\textwidth}
\begin{lstlisting}[basicstyle=\ttfamily\scriptsize]
/** LEGGE: la somma degli importi. */
public static double totale(Pagabile[] elenco) {
    double somma = 0;
    for (Pagabile p : elenco) {
        somma += p.importo();
    }
    return somma;
}

public static void main(String[] args) {
    Pagabile[] daPagare = {
        new Fattura("F-01", 3, 250),
        new Dipendente("Anna Russo", 1800),
        new Fattura("F-02", 1, 99.5)
    };
    for (Pagabile p : daPagare) {
        System.out.println(p + ": " + p.importo());
    }
    System.out.println(totale(daPagare));
\end{lstlisting}
\begin{lstlisting}[style=shell,basicstyle=\ttfamily\scriptsize\color{white}]
fattura F-01: 750.0
dipendente Anna Russo: 1800.0
fattura F-02: 99.5
2649.5
\end{lstlisting}

\column{0.44\textwidth}
\begin{block}{L'interfaccia è un tipo}
\small Variabili e array possono essere di tipo \texttt{Pagabile}, parametri compresi (variabili locali inizializzate dalla chiamata, lezione~3): un riferimento \texttt{Pagabile} porta a un oggetto di \textbf{qualunque} classe che la implementa (upcasting implicito, come con le superclassi).
\end{block}
\footnotesize
\begin{itemize}\itemsep=2pt
    \item \texttt{p.importo()}: \textbf{binding dinamico} (lezione~7): la JVM esegue l'\texttt{importo} della classe dell'oggetto. \texttt{totale} vale anche per classi future
    \item attraverso \texttt{p} si vedono solo i metodi di \texttt{Pagabile} (e di \texttt{Object}): \texttt{p.aumenta(100)} non compila. Per usarlo: \texttt{p instanceof Dipendente}, poi il cast \texttt{(Dipendente) p}. \texttt{instanceof} vale anche per le interfacce
    \item \texttt{new Pagabile()} non compila: gli oggetti nascono dalle classi
\end{itemize}
\end{columns}
\end{frame}

% ========================================================================
\begin{frame}[fragile]{Più interfacce per una classe -- versione 2}
\begin{columns}[T]
\column{0.58\textwidth}
\begin{lstlisting}[basicstyle=\ttfamily\scriptsize]
/** Un pagabile che si puo` pagare a rate. */
public interface Rateizzabile extends Pagabile {

    int numeroRate();  // public abstract
\end{lstlisting}
\begin{lstlisting}[basicstyle=\ttfamily\scriptsize]
/** Una fattura: a rate e stampabile. */
public class Fattura
        implements Rateizzabile, Stampabile {
\end{lstlisting}
\begin{lstlisting}[basicstyle=\ttfamily\scriptsize]
@Override
public int numeroRate() {
    return rate;
}

@Override
public void stampa() {
    System.out.println("FATTURA n. " + numero);
    System.out.println("totale " + importo()
            + " in " + rate + " rate");
}
\end{lstlisting}

\column{0.42\textwidth}
\MermaidAdjustBoxOpts{max width=\linewidth,max height=0.26\textheight,center}
\begin{mermaid}
---
config:
  themeVariables:
    fontSize: 28px
  class:
    hideEmptyMembersBox: true
---
classDiagram
  direction TB
  class Pagabile {
    <<interface>>
  }
  class Rateizzabile {
    <<interface>>
  }
  class Stampabile {
    <<interface>>
  }
  class Fattura
  class Dipendente
  Pagabile <|-- Rateizzabile
  Rateizzabile <|.. Fattura
  Stampabile <|.. Fattura
  Pagabile <|.. Dipendente
\end{mermaid}
\MermaidAdjustBoxOpts{max width=0.95\linewidth,max height=0.62\textheight,center}
\footnotesize
\begin{itemize}\itemsep=2pt
    \item \texttt{implements A, B}: la classe promette i metodi di \textbf{entrambe} (\texttt{Stampabile}: solo \texttt{stampa()}). Una sola superclasse, ma quanti tipi interfaccia vuole: \texttt{Fattura} è \texttt{Rateizzabile}, \texttt{Stampabile} e \texttt{Pagabile}: \emph{ereditarietà multipla} di tipi, senza stato
    \item \texttt{extends} e \texttt{implements} insieme: \texttt{class C extends S implements A, B}
    \item un'interfaccia può \textbf{estendere} un'altra: chi implementa \texttt{Rateizzabile} scrive \texttt{importo()} \textbf{e} \texttt{numeroRate()}
\end{itemize}
\end{columns}
\end{frame}

% ========================================================================
\begin{frame}[fragile]{Metodi \texttt{default} e \texttt{static} -- \texttt{ProvaPagabili.java} (versione 2)}
\begin{columns}[T]
\column{0.52\textwidth}
\begin{lstlisting}[basicstyle=\ttfamily\scriptsize]
public interface Pagabile {

    double importo();  // public abstract

    static double totale(Pagabile[] elenco) {
        double somma = 0;
        for (Pagabile p : elenco) {
            somma += p.importo();
        }
        return somma;
    }
}
\end{lstlisting}
\begin{lstlisting}[basicstyle=\ttfamily\scriptsize]
public interface Rateizzabile extends Pagabile {

    int numeroRate();  // public abstract

    default double rata() {
        return importo() / numeroRate();
    }
}
\end{lstlisting}

\column{0.48\textwidth}
\begin{lstlisting}[basicstyle=\ttfamily\scriptsize]
Fattura f = new Fattura("F-01", 3, 250, 2);
Dipendente d = new Dipendente("Anna", 1800);

Pagabile[] conti = { f, d };
System.out.println(Pagabile.totale(conti));
System.out.println(f.rata());   // default
// System.out.println(d.rata()); // errore
// System.out.println(f.totale()); // errore
\end{lstlisting}
\begin{lstlisting}[style=shell,basicstyle=\ttfamily\scriptsize\color{white}]
2550.0
375.0
\end{lstlisting}
\footnotesize
\begin{itemize}\itemsep=2pt
    \item un metodo \textbf{\texttt{default}} ha il corpo e lo \textbf{ereditano} le classi che implementano l'interfaccia: \texttt{f.rata()} funziona senza che \texttt{Fattura} l'abbia scritto. Usa solo i metodi dell'interfaccia (non c'è stato); si può ridefinire. \texttt{d.rata()} non compila: \texttt{d} non è \texttt{Rateizzabile}
    \item un metodo \textbf{\texttt{static}} si chiama con il nome dell'interfaccia: \texttt{Pagabile.totale(conti)}. \textbf{Diversamente dalle classi} (lezione~6), \texttt{f.totale()} \textbf{non compila}: i metodi \texttt{static} di un'interfaccia non passano alle classi che la implementano
\end{itemize}
\end{columns}
\end{frame}

% ========================================================================
\begin{frame}[fragile]{Interfacce o classi astratte? E finalmente \texttt{Cloneable}}
\begin{columns}[T]
\column{0.56\textwidth}
\footnotesize
\begin{center}
\renewcommand{\arraystretch}{1.0}
\scriptsize
\begin{tabular}{@{}>{\raggedright\arraybackslash}p{2.1cm}>{\raggedright\arraybackslash}p{2.8cm}>{\raggedright\arraybackslash}p{2.8cm}@{}}
\toprule
 & \textbf{classe astratta} (L7) & \textbf{interfaccia} (oggi) \\
\midrule
stato & variabili di istanza & no, solo costanti \\
costruttori & sì, con \texttt{super(...)} & no \\
metodi con il corpo & tutti quelli che servono & solo \texttt{default} e \texttt{static} \\
quante per classe & una (\texttt{extends}) & quante si vuole (\texttt{implements}) \\
la relazione & \emph{è-un}, con una parte comune da ereditare & \emph{sa fare}: un comportamento \\
\bottomrule
\end{tabular}
\end{center}
\begin{itemize}\itemsep=2pt
    \item \textbf{classe astratta} quando le sottoclassi condividono \emph{stato e codice} (\texttt{Forma}: \texttt{colore}, \texttt{piuGrandeDi}); \textbf{interfaccia} per classi non imparentate con gli stessi metodi, o per dare più tipi a una classe
    \item un ADT (lezione~5) è una struttura dati incapsulata dalle sue operazioni: la sua \emph{interfaccia} è l'elenco delle operazioni, il corpo la realizzazione. Un'\texttt{interface} di Java è proprio quell'elenco, separato dalle realizzazioni: nella lezione~12 la pila sarà un'\texttt{interface} realizzata da due classi, con un array e con una lista; il client non cambia
\end{itemize}

\column{0.44\textwidth}
\begin{lstlisting}[basicstyle=\ttfamily\scriptsize]
public interface Cloneable {
}
\end{lstlisting}
\begin{block}{\texttt{Cloneable}: nessun metodo}
\small Dalla lezione~8: \texttt{implements Cloneable} non obbliga a scrivere nulla, nemmeno \texttt{clone()}, che c'è già in \texttt{Object} (la classe lo ridefinisce \texttt{public} solo per offrirlo ai client). \textbf{Marca} la classe come copiabile: \texttt{Object.clone()} controlla se l'oggetto è \texttt{instanceof Cloneable}; se no, lancia \texttt{CloneNotSupportedException}: ecco il \texttt{throws} (lezione~10).
\end{block}
\footnotesize
\begin{itemize}\itemsep=2pt
    \item \texttt{Cloneable} sta in \texttt{java.lang}, come \texttt{Object}: niente \texttt{import}
    \item lezione~12: \texttt{Comparable}, con \texttt{compareTo}, per \texttt{Arrays.sort} sui nostri oggetti
\end{itemize}
\end{columns}
\end{frame}

% ========================================================================
\section{I package}

\begin{frame}[fragile]{Package: il concetto -- \texttt{contabilita\_pkg/}}
\begin{columns}[T]
\column{0.58\textwidth}
\begin{block}{Package}
\small Un \emph{package} è un gruppo di classi e interfacce correlate, con un nome: è l'unità con cui Java organizza il codice in cartelle e librerie. \texttt{java.util} è il package con \texttt{Scanner}, \texttt{Arrays} e \texttt{Objects}; \texttt{contabilita} è il nostro.
\end{block}
\footnotesize
\begin{itemize}\itemsep=2pt
    \item si dichiara con \texttt{package contabilita;} come \textbf{prima riga} del file (prima solo commenti); il file sta nella \textbf{cartella con lo stesso nome} del package: \texttt{contabilita/Fattura.java}
    \item il \textbf{nome qualificato} di una classe è package più classe: \texttt{contabilita.Fattura}; due classi \texttt{Persona} in package diversi convivono
    \item i package si \textbf{innestano}: \texttt{java.util.Arrays} è la classe \texttt{Arrays} nel package \texttt{util} dentro \texttt{java}. Nomi in minuscolo
    \item finora nessun file aveva la riga \texttt{package}: le classi senza package stanno nel \textbf{package di default}, senza nome, che \texttt{javac} e \texttt{java} cercano nella cartella corrente. Un package raggruppa classi coese: il \emph{modulo} della lezione~4, su scala maggiore
\end{itemize}

\column{0.42\textwidth}
\MermaidAdjustBoxOpts{max width=\linewidth,max height=0.6\textheight,center}
\begin{mermaid}
---
config:
  themeVariables:
    fontSize: 20px
  flowchart:
    nodeSpacing: 10
    rankSpacing: 26
    wrappingWidth: 600
---
flowchart TB
  R["<b>contabilita_pkg/</b><br/>la radice: qui si compila"]
  C["<b>contabilita/</b><br/>package contabilita"]
  A["<b>app/</b><br/>package app"]
  C1["Pagabile.java<br/>Rateizzabile.java<br/>Stampabile.java<br/>Fattura.java<br/>Dipendente.java"]
  A1["ProvaPagabili.java"]
  R --> C --> C1
  R --> A --> A1
  style R fill:#fef3c7,stroke:#b45309
  style C fill:#dbeafe,stroke:#1d4ed8
  style A fill:#dcfce7,stroke:#15803d
  style C1 fill:#dbeafe,stroke:#1d4ed8
  style A1 fill:#dcfce7,stroke:#15803d
\end{mermaid}
\MermaidAdjustBoxOpts{max width=0.95\linewidth,max height=0.62\textheight,center}
\end{columns}
\end{frame}

% ========================================================================
\begin{frame}[fragile]{Compilare ed eseguire con i package}
\begin{columns}[T]
\column{0.5\textwidth}
{\footnotesize In \texttt{contabilita/Fattura.java}, prima riga:}
\begin{lstlisting}[basicstyle=\ttfamily\scriptsize]
package contabilita;

/** Una fattura: a rate e stampabile. */
public class Fattura
        implements Rateizzabile, Stampabile {
\end{lstlisting}
{\footnotesize Dal terminale, nella cartella \textbf{radice} \texttt{contabilita\_pkg/}:}
\begin{lstlisting}[style=shell,basicstyle=\ttfamily\scriptsize\color{white}]
$ javac app/ProvaPagabili.java
$ java app.ProvaPagabili
2550.0
FATTURA n. F-01
totale 750.0 in 2 rate
\end{lstlisting}
{\footnotesize Dentro \texttt{app/}, invece, \textbf{non} funziona:}
\begin{lstlisting}[style=shell,basicstyle=\ttfamily\scriptsize\color{white}]
$ cd app
$ java ProvaPagabili
Error: Could not find or load main class
ProvaPagabili
Caused by: java.lang.NoClassDefFoundError:
ProvaPagabili (wrong name: app/ProvaPagabili)
\end{lstlisting}

\column{0.5\textwidth}
\footnotesize
\begin{itemize}\itemsep=3pt
    \item ci si mette nella cartella \textbf{superiore} ai package, la radice del progetto
    \item \texttt{javac} riceve il \textbf{percorso} del file con il \texttt{main}. Come nella lezione~4, compila da solo anche le classi che servono, cercandole dalla cartella corrente: \texttt{contabilita.Fattura} in \texttt{contabilita/Fattura.java}. I \texttt{.class} finiscono accanto ai sorgenti
    \item \texttt{java} riceve il \textbf{nome qualificato} della classe: non un percorso, non \texttt{.class}. La JVM lo traduce in \texttt{app/ProvaPagabili.class} dalla cartella corrente: per questo da dentro \texttt{app/} non lo trova (\emph{wrong name})
    \item finora tutte le nostre classi stavano nel package di default, lo stesso per tutte, e si vedevano con il nome semplice. Una classe del package di default, però, \textbf{non ha nome qualificato}: da \texttt{app} non si può né importare né nominare. Quindi, quando un progetto adotta i package, \textbf{tutte} le classi vanno in un package, anche quella con il \texttt{main}
    \item in VS Code basta aprire la cartella radice e usare \emph{Run}
\end{itemize}
\end{columns}
\end{frame}

% ========================================================================
\begin{frame}[fragile]{\texttt{import}: usare le classi di un altro package -- \texttt{app/ProvaPagabili.java}}
\begin{columns}[T]
\column{0.57\textwidth}
\begin{lstlisting}[basicstyle=\ttfamily\scriptsize]
package app;

import contabilita.Fattura;
import contabilita.Pagabile;

/** Il client, in un altro package. */
public class ProvaPagabili {

    public static void main(String[] args) {
        Fattura f = new Fattura("F-01", 3, 250, 2);
        contabilita.Dipendente d;   // qualificato
        d = new contabilita.Dipendente("Anna", 1800);
        Pagabile[] conti = { f, d };
        System.out.println(Pagabile.totale(conti));
        f.stampa();
    }
}
\end{lstlisting}
\footnotesize
\texttt{Dipendente} non è importato: si usa con il nome qualificato. \texttt{Fattura}, invece, usa \texttt{Rateizzabile} e \texttt{Stampabile} senza \texttt{import}: le classi dello \textbf{stesso} package si vedono con il nome semplice.

\column{0.43\textwidth}
\begin{block}{\texttt{import}}
\small Fuori dal suo package una classe si chiama con il nome qualificato; \texttt{import contabilita.Fattura;}, in testa al file, permette di scrivere solo \texttt{Fattura}. \texttt{import} non copia nulla: dice al compilatore \textbf{dove cercare}.
\end{block}
\footnotesize
\begin{itemize}\itemsep=2pt
    \item tre forme: nome qualificato; \texttt{import} di una classe; \texttt{import contabilita.*;} per tutto il package
    \item classi con lo \textbf{stesso nome} in package diversi: se ne importa una sola; se servono entrambe, i due \texttt{import} sono un errore: la seconda si usa con il nome qualificato
    \item \texttt{java.lang} è importato \textbf{da solo}: \texttt{String}, \texttt{System}, \texttt{Math}, \texttt{Object}. \texttt{java.util} no: \texttt{Scanner}, \texttt{Arrays} e \texttt{Objects} li abbiamo sempre importati
\end{itemize}
\end{columns}
\end{frame}

% ========================================================================
\begin{frame}[fragile]{Visibilità di package e riepilogo dei modificatori}
\begin{columns}[T]
\column{0.5\textwidth}
\footnotesize
\begin{center}
\renewcommand{\arraystretch}{1.25}
\scriptsize
\begin{tabular}{@{}lcccc@{}}
\toprule
\textbf{membro} & \textbf{stessa} & \textbf{stesso} & \textbf{sottoclasse in} & \textbf{ovunque} \\
 & \textbf{classe} & \textbf{package} & \textbf{altro package} & \\
\midrule
\texttt{private} & sì & -- & -- & -- \\
(nessuno) & sì & sì & -- & -- \\
\texttt{protected} & sì & sì & sì & -- \\
\texttt{public} & sì & sì & sì & sì \\
\bottomrule
\end{tabular}
\end{center}
\begin{block}{Visibilità di package}
\small Un membro \textbf{senza modificatore} è visibile da tutte le classi del suo package e da nessun'altra. Una \textbf{classe senza \texttt{public}} è visibile solo nel suo package: fuori non si può nominare, né importare. Il file non conta: conta il package.
\end{block}

\column{0.5\textwidth}
\footnotesize
\begin{itemize}\itemsep=3pt
    \item il \texttt{protected} della lezione~7 è ora completo: sottoclassi \textbf{e} tutto il package
    \item per le classi solo due scelte: \texttt{public} (al più una per file, con il nome del file) o niente (visibile nel package): una classe scritta in un file non può essere \texttt{private}
    \item nelle interfacce ciò che si offre è \texttt{public}: costanti, metodi, e i metodi che le implementano
    \item la regola di sempre, completata: \texttt{private} per lo stato; \texttt{public} per l'interfaccia offerta ai client; \textbf{niente modificatore} per le classi e i metodi di servizio che devono vedere solo le classi dello stesso package; \texttt{protected} raramente
\end{itemize}
\end{columns}
\end{frame}

% ========================================================================
\section{Esercitazione}

\begin{frame}[fragile]{Commentiamo insieme (20 min) -- il problema}
\footnotesize
\textbf{Il problema.} Un simulatore di casa domotica gestisce dispositivi diversi: lampade, ventilatori, un allarme. Il programma deve poterli \textbf{accendere e spegnere tutti insieme}, senza sapere di che tipo è ciascuno, e \textbf{regolare} quelli che si regolano: la luce di una lampada, la velocità di un ventilatore. L'allarme è solo acceso o spento. Le due interfacce, le tre classi e il \texttt{Simulatore} sono già scritti: li leggiamo riga per riga, cercando che cosa promette ogni interfaccia, chi implementa che cosa e dove serve \texttt{instanceof}.

\vspace{0.05cm}
\begin{columns}[T]
\column{0.5\textwidth}
\textbf{Le regole}
\begin{itemize}\itemsep=2pt
    \item ogni dispositivo si accende, si spegne, dice se è acceso e si può \emph{commutare} (spegnere se acceso, accendere se spento)
    \item i dispositivi regolabili hanno un \textbf{livello} che si alza e si abbassa di un passo alla volta, solo se accesi, entro i limiti: la lampada da 1 a 10 (parte da 5), il ventilatore da 1 a 3 (parte da 1)
    \item l'allarme non è regolabile
    \item il programma stampa ogni dispositivo e quanti sono accesi
\end{itemize}

\column{0.5\textwidth}
\MermaidAdjustBoxOpts{max width=\linewidth,max height=0.5\textheight,center}
\begin{mermaid}
---
config:
  class:
    hideEmptyMembersBox: true
---
classDiagram
  direction TB
  class Accendibile {
    <<interface>>
    +accendi()
    +spegni()
    +isAcceso() boolean
    +commuta()
    +quantiAccesi(Accendibile[])$ int
  }
  class Regolabile {
    <<interface>>
    +alza()
    +abbassa()
    +getLivello() int
  }
  class Lampada
  class Ventilatore
  class Allarme
  Accendibile <|.. Lampada
  Regolabile <|.. Lampada
  Accendibile <|.. Ventilatore
  Regolabile <|.. Ventilatore
  Accendibile <|.. Allarme
\end{mermaid}
\MermaidAdjustBoxOpts{max width=0.95\linewidth,max height=0.62\textheight,center}
\end{columns}
\end{frame}

% ========================================================================
\begin{frame}[fragile]{Commentiamo insieme -- \texttt{Accendibile.java}}
\begin{columns}[T]
\column{0.48\textwidth}
\begin{lstlisting}[basicstyle=\ttfamily\scriptsize]
/** Un dispositivo: si accende e spegne. */
public interface Accendibile {

    void accendi();

    void spegni();

    boolean isAcceso();

    /** default: inverte lo stato. */
    default void commuta() {
        if (isAcceso()) {
            spegni();
        } else {
            accendi();
        }
    }
\end{lstlisting}

\column{0.52\textwidth}
\begin{lstlisting}[basicstyle=\ttfamily\scriptsize]
/** static: LEGGE quanti sono accesi. */
static int quantiAccesi(Accendibile[] elenco) {
    int n = 0;
    for (Accendibile d : elenco) {
        if (d.isAcceso()) {
            n++;
        }
    }
    return n;
}
\end{lstlisting}
\footnotesize
\begin{itemize}\itemsep=2pt
    \item tre metodi astratti: ogni classe che implementa \texttt{Accendibile} li scriverà a modo suo
    \item \texttt{commuta} è \texttt{default}: scritto una volta sola in termini di \texttt{isAcceso}, \texttt{spegni} e \texttt{accendi}, che l'interfaccia non sa come sono fatti ma sa che ci saranno. Lo ereditano lampade, ventilatori e allarme
    \item \texttt{quantiAccesi} è \texttt{static}: lavora su un array di \texttt{Accendibile} qualsiasi e si chiama \texttt{Accendibile.quantiAccesi(...)}
\end{itemize}
\end{columns}
\end{frame}

% ========================================================================
\begin{frame}[fragile]{Commentiamo insieme -- \texttt{Regolabile.java} e \texttt{Lampada.java} (1/2)}
\begin{columns}[T]
\column{0.5\textwidth}
\begin{lstlisting}[basicstyle=\ttfamily\scriptsize]
/** Ha un livello regolabile a passi. */
public interface Regolabile {

    void alza();

    void abbassa();

    int getLivello();
}
\end{lstlisting}
\footnotesize
\begin{itemize}\itemsep=2pt
    \item le due interfacce sono indipendenti: un dispositivo può implementarne una sola o entrambe
    \item \texttt{Lampada} e \texttt{Ventilatore} le implementano entrambe, ma \textbf{non} sono imparentate: nessuna superclasse comune, ciascuna con le sue variabili di istanza (\texttt{luce}, \texttt{velocita}). Hanno in comune solo i due tipi interfaccia
    \item \texttt{Ventilatore} è fatto come \texttt{Lampada}: cambiano i nomi e i limiti (\texttt{velocita} da 1 a 3)
\end{itemize}

\column{0.5\textwidth}
\begin{lstlisting}[basicstyle=\ttfamily\scriptsize]
/** Una lampada: luce regolabile da 1 a 10. */
public class Lampada
        implements Accendibile, Regolabile {

    private static final int MAX = 10;

    private final String stanza;
    private boolean accesa;
    private int luce;

    public Lampada(String stanza) {
        this.stanza = stanza;
        accesa = false;
        luce = 5;
    }

    @Override
    public void accendi() {
        accesa = true;
    }
\end{lstlisting}
\end{columns}
\end{frame}

% ========================================================================
\begin{frame}[fragile]{Commentiamo insieme -- \texttt{Lampada.java} (2/2) e \texttt{Allarme.java}}
\begin{columns}[T]
\column{0.5\textwidth}
\begin{lstlisting}[basicstyle=\ttfamily\scriptsize]
    @Override
    public void spegni() {
        accesa = false;
    }

    @Override
    public boolean isAcceso() {
        return accesa;
    }

    /** MODIFICA: piu` luce, se accesa. */
    @Override
    public void alza() {
        if (accesa && luce < MAX) {
            luce++;
        }
    }
\end{lstlisting}
\footnotesize
I sei metodi promessi ci sono tutti, \texttt{public} e con \texttt{@Override}. Non mostrati: \texttt{abbassa()}, speculare ad \texttt{alza()}; \texttt{getLivello()}, che restituisce \texttt{luce}; \texttt{toString()}: \texttt{lampada cucina: accesa, luce 8/10} oppure \texttt{lampada cucina: spenta}.

\column{0.5\textwidth}
\begin{lstlisting}[basicstyle=\ttfamily\scriptsize]
/** Un allarme: acceso o spento. */
class Allarme implements Accendibile {

    private boolean attivo;
\end{lstlisting}
\begin{lstlisting}[basicstyle=\ttfamily\scriptsize]
    @Override
    public void accendi() {
        attivo = true;
    }

    @Override
    public void spegni() {
        attivo = false;
    }

    @Override
    public boolean isAcceso() {
        return attivo;
    }
\end{lstlisting}
\footnotesize
\texttt{Allarme} è dichiarata \textbf{senza} \texttt{public}: finché tutto sta nella stessa cartella non cambia nulla (ci torniamo fra poco). Non mostrati: il costruttore, che mette \texttt{attivo} a \texttt{false}, e \texttt{toString} (\texttt{allarme: attivo} o \texttt{allarme: disattivo}).
\end{columns}
\end{frame}

% ========================================================================
\begin{frame}[fragile]{Commentiamo insieme -- \texttt{Simulatore.java}}
\begin{columns}[T]
\column{0.52\textwidth}
\begin{lstlisting}[basicstyle=\ttfamily\scriptsize]
Accendibile[] casa = {
    new Lampada("cucina"),
    new Ventilatore("studio"),
    new Allarme(),
    new Lampada("camera")
};
for (Accendibile d : casa) {
    d.accendi();
}
for (Accendibile d : casa) {
    if (d instanceof Regolabile) {
        Regolabile r = (Regolabile) d;
        r.alza();
        r.alza();
        r.alza();
    }
}
casa[2].commuta();      // allarme: si spegne
casa[3].spegni();
for (Accendibile d : casa) {
    System.out.println(d);
}
System.out.println("accesi: "
        + Accendibile.quantiAccesi(casa));
\end{lstlisting}

\column{0.48\textwidth}
\begin{lstlisting}[style=shell,basicstyle=\ttfamily\scriptsize\color{white}]
lampada cucina: accesa, luce 8/10
ventilatore studio: acceso, velocita 3/3
allarme: disattivo
lampada camera: spenta
accesi: 2
\end{lstlisting}
\footnotesize
\textbf{Rispondiamo insieme:}
\begin{enumerate}\itemsep=1pt
    \item perché \texttt{casa} può contenere lampade, ventilatori e allarmi, che non hanno una superclasse comune?
    \item nel secondo ciclo, perché \texttt{d.alza()} non compilerebbe? E a che cosa servono \texttt{instanceof} e il cast?
    \item il ventilatore riceve tre \texttt{alza()} e arriva a 3, non a 4; la lampada a 8: dove sta la differenza?
    \item \texttt{Allarme} non ha scritto \texttt{commuta}: perché \texttt{casa[2].commuta()} funziona?
    \item se togliessimo \texttt{isAcceso} da \texttt{Lampada}, quale errore darebbe il compilatore, e su quale classe?
\end{enumerate}
\end{columns}
\end{frame}

% ========================================================================
\begin{frame}[fragile]{Insieme, passo passo (30 min) -- il problema}
\begin{columns}[T]
\column{0.58\textwidth}
\footnotesize
\textbf{Il problema.} Il simulatore cresce: altri dispositivi, altri programmi che li usano. Le due interfacce e le tre classi dei dispositivi sono \textbf{riusabili}, il \texttt{Simulatore} no: è uno dei tanti programmi possibili. Vogliamo \textbf{riorganizzare il progetto in package}: i dispositivi formano una piccola libreria, separata dai programmi che la usano.

\vspace{0.1cm}
\textbf{Le regole}
\begin{itemize}\itemsep=2pt
    \item il codice delle classi non cambia: cambiano la posizione dei file e le righe iniziali
    \item un package \texttt{dispositivi} per le interfacce e le classi dei dispositivi, un package \texttt{app} per il \texttt{Simulatore}
    \item tutto ciò che \texttt{app} usa deve essere \textbf{visibile} da un altro package
    \item alla fine il programma si compila e si esegue dalla radice e stampa quello di prima
\end{itemize}

\vspace{0.1cm}
Partiamo dalla cartella \texttt{simulatore/} con i sei file. Un passo alla volta, compilando dal terminale per vedere \textbf{gli errori} che ogni passo produce finché il successivo non lo completa.

\column{0.42\textwidth}
\MermaidAdjustBoxOpts{max width=\linewidth,max height=0.7\textheight,center}
\begin{mermaid}
---
config:
  themeVariables:
    fontSize: 20px
  flowchart:
    nodeSpacing: 10
    rankSpacing: 26
    wrappingWidth: 600
---
flowchart TB
  S["<b>simulatore/</b><br/>Accendibile.java<br/>Regolabile.java<br/>Lampada.java<br/>Ventilatore.java<br/>Allarme.java<br/>Simulatore.java"]
  R["<b>simulatore_pkg/</b>"]
  D["<b>dispositivi/</b><br/>Accendibile.java<br/>Regolabile.java<br/>Lampada.java<br/>Ventilatore.java<br/>Allarme.java"]
  A["<b>app/</b><br/>Simulatore.java"]
  S -- "riorganizziamo" --> R
  R --> D
  R --> A
  style S fill:#fef3c7,stroke:#b45309
  style R fill:#fef3c7,stroke:#b45309
  style D fill:#dbeafe,stroke:#1d4ed8
  style A fill:#dcfce7,stroke:#15803d
\end{mermaid}
\MermaidAdjustBoxOpts{max width=0.95\linewidth,max height=0.62\textheight,center}
\end{columns}
\end{frame}

% ========================================================================
\begin{frame}[fragile]{Insieme, passo passo -- le decisioni}
\begin{columns}[T]
\column{0.6\textwidth}
\footnotesize
\textbf{Prima di toccare i file, decidiamo insieme:}
\begin{enumerate}\itemsep=3pt
    \item \textbf{quali package?} \textrightarrow{} \texttt{dispositivi}: le due interfacce e le tre classi, coese fra loro e senza dipendenze dal programma; \texttt{app}: il solo \texttt{Simulatore}, che dipende da \texttt{dispositivi} e non viceversa
    \item \textbf{le cartelle?} \textrightarrow{} una radice \texttt{simulatore\_pkg/} con dentro \texttt{dispositivi/} e \texttt{app/}: stesso nome dei package, lettera per lettera
    \item \textbf{la prima riga di ogni file?} \textrightarrow{} \texttt{package dispositivi;} nei cinque file dei dispositivi, \texttt{package app;} in \texttt{Simulatore.java}
    \item \textbf{\texttt{Simulatore} come trova le classi?} \textrightarrow{} con gli \texttt{import}: uno per ogni nome che usa (\texttt{Accendibile}, \texttt{Regolabile}, \texttt{Lampada}, \texttt{Ventilatore}, \texttt{Allarme}), oppure \texttt{import dispositivi.*;}
    \item \textbf{che cosa deve essere \texttt{public}?} \textrightarrow{} tutto ciò che \texttt{app} nomina: le due interfacce e le tre classi, con i loro costruttori. Le interfacce e \texttt{Lampada}, \texttt{Ventilatore} lo sono già; \texttt{Allarme} \textbf{no}: il compilatore ce lo dirà
    \item \textbf{da dove si compila?} \textrightarrow{} dalla radice: \texttt{javac app/Simulatore.java}, poi \texttt{java app.Simulatore}
\end{enumerate}

\column{0.4\textwidth}
\footnotesize
\textbf{Gli errori che si incontrano}
\begin{itemize}\itemsep=3pt
    \item file spostati in \texttt{dispositivi/}, \texttt{Simulatore} ancora senza \texttt{import}: \emph{``cannot find symbol: class Lampada''}
    \item \texttt{import} messi, ma in un file manca la riga \texttt{package}: \emph{``cannot access Accendibile: bad source file \ldots{} file does not contain class dispositivi.Accendibile''}
    \item \texttt{import dispositivi.Allarme;} con \texttt{Allarme} senza \texttt{public}: \emph{``Allarme is not public in dispositivi; cannot be accessed from outside package''}
    \item \texttt{java Simulatore} da dentro \texttt{app/}: \emph{``Could not find or load main class Simulatore''}
\end{itemize}
\vspace{0.1cm}
Ogni messaggio dice esattamente qual è il passo che manca.
\end{columns}
\end{frame}

% ========================================================================
\begin{frame}[fragile]{Una possibile soluzione (1/2) -- il package \texttt{dispositivi}}
\begin{columns}[T]
\column{0.5\textwidth}
{\footnotesize \texttt{dispositivi/Accendibile.java}, prime righe:}
\begin{lstlisting}[basicstyle=\ttfamily\scriptsize]
package dispositivi;

/** Un dispositivo: si accende e spegne. */
public interface Accendibile {
\end{lstlisting}
{\footnotesize \texttt{dispositivi/Regolabile.java}:}
\begin{lstlisting}[basicstyle=\ttfamily\scriptsize]
package dispositivi;

/** Ha un livello regolabile a passi. */
public interface Regolabile {
\end{lstlisting}
{\footnotesize \texttt{dispositivi/Lampada.java} (e così \texttt{Ventilatore.java}):}
\begin{lstlisting}[basicstyle=\ttfamily\scriptsize]
package dispositivi;

/** Una lampada: luce regolabile da 1 a 10. */
public class Lampada
        implements Accendibile, Regolabile {
\end{lstlisting}

\column{0.5\textwidth}
{\footnotesize \texttt{dispositivi/Allarme.java}, ora \texttt{public}:}
\begin{lstlisting}[basicstyle=\ttfamily\scriptsize]
package dispositivi;

/** Un allarme: acceso o spento. */
public class Allarme implements Accendibile {
\end{lstlisting}
\footnotesize
\begin{itemize}\itemsep=3pt
    \item il resto dei cinque file è \textbf{identico} alla versione di prima: il codice non sa in che package sta
    \item \texttt{Lampada} usa \texttt{Accendibile} e \texttt{Regolabile} senza \texttt{import}: stesso package
    \item \texttt{Allarme} senza \texttt{public} compilava da sola e con le altre: \texttt{dispositivi} la vede. È \texttt{app} che non la vedeva: la visibilità di package, finora invisibile, qui conta. Il suo costruttore era già \texttt{public}: anche lui deve esserlo, altrimenti \texttt{new Allarme()} da \texttt{app} fallisce con un messaggio analogo, sul costruttore
    \item nessun \texttt{import} di \texttt{java.lang}: \texttt{String} e \texttt{System} si vedono sempre
\end{itemize}
\end{columns}
\end{frame}

% ========================================================================
\begin{frame}[fragile]{Una possibile soluzione (2/2) -- il package \texttt{app}}
\begin{columns}[T]
\column{0.5\textwidth}
{\footnotesize \texttt{app/Simulatore.java}, prime righe:}
\begin{lstlisting}[basicstyle=\ttfamily\scriptsize]
package app;

import dispositivi.Accendibile;
import dispositivi.Regolabile;
import dispositivi.Lampada;
import dispositivi.Ventilatore;
import dispositivi.Allarme;

/** Dispositivi diversi, tutti Accendibile. */
public class Simulatore {

    public static void main(String[] args) {
        Accendibile[] casa = {
            new Lampada("cucina"),
\end{lstlisting}
\footnotesize
Il \texttt{main} prosegue come prima, senza una modifica. Cinque \texttt{import} o uno solo, \texttt{import dispositivi.*;}: con i nomi espliciti si legge subito da dove viene ogni classe.

\column{0.5\textwidth}
{\footnotesize Dalla radice \texttt{simulatore\_pkg/}:}
\begin{lstlisting}[style=shell,basicstyle=\ttfamily\scriptsize\color{white}]
$ javac app/Simulatore.java
$ java app.Simulatore
lampada cucina: accesa, luce 8/10
ventilatore studio: acceso, velocita 3/3
allarme: disattivo
lampada camera: spenta
accesi: 2
\end{lstlisting}
\footnotesize
\begin{itemize}\itemsep=3pt
    \item \texttt{javac} compila \texttt{Simulatore} e, seguendo gli \texttt{import}, i cinque file di \texttt{dispositivi/}: i \texttt{.class} compaiono nelle due cartelle
    \item \texttt{java app.Simulatore}: nome qualificato, dalla radice
    \item stesso output di prima: la riorganizzazione \textbf{non} cambia il comportamento, cambia la struttura. Un nuovo programma in \texttt{app/} riusa i dispositivi con gli stessi \texttt{import}; un nuovo dispositivo va in \texttt{dispositivi/}, \texttt{public}, con la sua riga \texttt{package}: è l'homework
\end{itemize}
\end{columns}
\end{frame}

% ========================================================================
\section{Wrap-up}

\begin{frame}{Riepilogo: cosa abbiamo trattato}
\footnotesize
\begin{center}
\renewcommand{\arraystretch}{1.1}
\scriptsize
\begin{tabular}{|>{\raggedright\arraybackslash}p{6.2cm}|>{\raggedright\arraybackslash}p{7.6cm}|}
\hline
\textbf{Concetto (unità)} & \textbf{Dove compare oggi} \\
\hline
\texttt{interface}: solo intestazioni e costanti; quando usarla (U17) & \texttt{Pagabile}; \texttt{Accendibile}, \texttt{Regolabile} \\
\hline
\texttt{implements}, metodi \texttt{public}, errore se ne manca uno (U17) & \texttt{Fattura}, \texttt{Dipendente}; \texttt{Lampada}, \texttt{Ventilatore}, \texttt{Allarme} \\
\hline
L'interfaccia come tipo: polimorfismo, \texttt{instanceof}, niente \texttt{new} (U17) & \texttt{Pagabile[] daPagare}; \texttt{Accendibile[] casa}, \texttt{instanceof Regolabile} \\
\hline
Più interfacce per una classe; un'interfaccia ne estende un'altra (U17) & \texttt{Fattura implements Rateizzabile, Stampabile}; \texttt{Rateizzabile extends Pagabile} \\
\hline
Metodi \texttt{default} e \texttt{static} (U17) & \texttt{rata()}, \texttt{Pagabile.totale}; \texttt{commuta()}, \texttt{Accendibile.quantiAccesi} \\
\hline
Interfacce o classi astratte; \texttt{Cloneable} (U17) & la tabella; \texttt{implements Cloneable} della lezione~8 \\
\hline
\texttt{package}, nome qualificato, cartelle, package di default (U18) & \texttt{contabilita\_pkg}; \texttt{simulatore\_pkg} \\
\hline
Compilare ed eseguire dalla radice; \texttt{import}; \texttt{java.lang} (U18) & \texttt{javac app/Simulatore.java}; \texttt{java app.Simulatore}; \texttt{app/ProvaPagabili} \\
\hline
Visibilità di package e riepilogo dei modificatori (U18) & la tabella; \texttt{Allarme} senza \texttt{public} \\
\hline
\end{tabular}
\end{center}

\vspace{0.1cm}
\textbf{Prossima lezione (L10)}: le eccezioni -- \texttt{try}/\texttt{catch}/\texttt{finally}, lo stack trace, \texttt{throws} (finalmente \texttt{CloneNotSupportedException}), eccezioni controllate e non, eccezioni nostre; la ricorsione -- casi base e ricorsivi, lo stack delle chiamate, da ricorsivo a iterativo.
\end{frame}

% ========================================================================
\begin{frame}{Homework 9 (da svolgere entro la prossima lezione)}
\footnotesize\setlength{\parskip}{0pt}\vspace{-0.15cm}
Tre esercizi, ciascuno con il suo programma di prova; il secondo e il terzo usano i package.
\begin{enumerate}\itemsep=1pt
    \item \textbf{\texttt{Spedibile.java}} (con \texttt{Pacco.java} e \texttt{Lettera.java}): un'interfaccia con la costante \texttt{TARIFFA\_AL\_KG} = 2.0, il metodo \texttt{peso()} (in kg) e il metodo \texttt{default} \texttt{costoSpedizione()} = 3.0 + tariffa $\times$ peso. \texttt{Pacco} (destinatario, peso) usa il costo \texttt{default}; \texttt{Lettera} (destinatario, grammi \texttt{int}) converte in kg e \textbf{ridefinisce} \texttt{costoSpedizione()}: 1.20 fisso. \texttt{ProvaSpedizioni}: uno \texttt{Spedibile[]} con due pacchi e due lettere; stampa il costo di ciascuno (\texttt{printf}, due decimali), il totale con un suo metodo \texttt{static} \texttt{costoTotale(Spedibile[])}, quante lettere ci sono (\texttt{instanceof}), la costante.
    \item \textbf{\texttt{Termostato.java}}: copiate la cartella \texttt{simulatore\_pkg} di oggi in \texttt{homework/termostato/}. Nel package \texttt{dispositivi} aggiungete \texttt{Termostato}, che implementa \texttt{Accendibile} e \texttt{Regolabile}: acceso o spento, temperatura \texttt{int} da 15 a 30 gradi (parte da 20), \texttt{alza}/\texttt{abbassa} di un grado solo se acceso; \texttt{toString} come \texttt{termostato: acceso, 22 gradi} (o \texttt{spento, 20 gradi}). In \texttt{app}, \texttt{ProvaTermostato}: un \texttt{Accendibile[]} con un termostato e una lampada; alza il termostato da spento (non cambia), accende tutti, alza due volte i regolabili con \texttt{instanceof}, stampa il termostato e il suo livello, lo commuta e ristampa, conta gli accesi con \texttt{quantiAccesi}.
    \item \textbf{\texttt{Magazzino.java}} (package \texttt{magazzino}, con \texttt{Articolo.java} e \texttt{Scorta.java}): \texttt{Articolo} è \texttt{public} e immutabile (codice, nome, prezzo; \texttt{toString}; \texttt{equals} e \texttt{hashCode} sul solo codice). \texttt{Scorta} è \textbf{senza} \texttt{public}, come il costruttore e i metodi (tranne \texttt{toString}); variabili \texttt{private}: un articolo e la quantità; \texttt{aggiungi(int n)}, \texttt{togli(int n)} (\texttt{false} se i pezzi non bastano), \texttt{valore()}. \texttt{Magazzino} è \texttt{public}: uno \texttt{Scorta[]} con la capienza data al costruttore; \texttt{carica(Articolo a, int n)} aggiunge alla scorta di \texttt{a} se c'è (cercata con \texttt{equals}), altrimenti la crea (\texttt{false} se non c'è posto); \texttt{scarica(Articolo a, int n)}; \texttt{disponibili(Articolo a)}; \texttt{valoreTotale()}; \texttt{toString}. \texttt{app.ProvaMagazzino} carica due articoli, ricarica il primo con un \texttt{Articolo} nuovo con lo stesso codice, scarica più del disponibile e poi meno, stampa disponibilità, valore e magazzino; prova \texttt{new magazzino.Scorta(...)}: leggete l'errore, poi commentate la riga.
\end{enumerate}
\end{frame}

% ========================================================================
\begin{frame}
\begin{center}
{\Huge Domande?}\\[2em]
{\normalsize \href{mailto:mattia.fonisto@unina.it}{\texttt{mattia.fonisto@unina.it}}}
\end{center}
\end{frame}

\end{document}
